% Partie : théorie de la démonstration
% Chapitre : propositions comme types
% Section : règles de tN2

\documentclass[../../../include/open-logic-section]{subfiles}

\begin{document}

\begin{table}
\begin{tabular}{@{}c@{}c@{}}
\multicolumn{2}{l}{\textbf{Axiomes :}}\\
\multicolumn{2}{c}{$\Gamma \Sequent x:!A$ si $x:!A \in \Gamma$}\\[3ex]
\multicolumn{2}{l}{\textbf{Règles d'inférence :}}\\[3ex]
\multicolumn{2}{c}{\Axiom$x: !A, \Gamma \fCenter N:!B$
\RightLabel{\Intro{\lif}}
\UnaryInf$\Gamma \fCenter \lambd[\typeof{x}{!A}][N] : !A \lif !B $
\DisplayProof}
\\[4ex]
\multicolumn{2}{c}{\Axiom$\Gamma_1 \fCenter N:!A \lif !B$
\Axiom$\Gamma_2 \fCenter M:!A$
\RightLabel{\Elim{\lif}}
\BinaryInf$\Gamma_1, \Gamma_2 \fCenter (NM) : !B$
\DisplayProof}
\\[4ex]
\multicolumn{2}{c}{\Axiom$\Gamma_1 \fCenter N_1 : !A$
\Axiom$\Gamma_2 \fCenter N_2 : !B$
\RightLabel{\Intro{\land}}
\BinaryInf$\Gamma_1, \Gamma_2 \fCenter \pair{N_1}{N_2} : !A \land !B $
\DisplayProof}
\\[4ex]
\Axiom$\Gamma \fCenter N: !A \land !B$
\RightLabel{\Elim{\land}[1]}
\UnaryInf$\Gamma \fCenter \proj{1}{N} : !A$
\DisplayProof
&
\Axiom$\Gamma \fCenter N:!A \land !B$
\RightLabel{\Elim{\land}[2]}
\UnaryInf$\Gamma \fCenter \proj{2}{N} : !B$
\DisplayProof\\[3ex]
\\[4ex]
\Axiom$\Gamma \fCenter N: !A$
\RightLabel{\Intro{\lor}[1]}
\UnaryInf$\Gamma \fCenter \inj[!B]{1}{N} : !A \lor !B$
\DisplayProof
&
\Axiom$ \Gamma \fCenter N : !B$
\RightLabel{\Intro{\lor}[2]}
\UnaryInf$\Gamma \fCenter \inj[!A]{2}{N} : !A \lor !B$
\DisplayProof\\[3ex]
\multicolumn{2}{c}{
\begin{tabular}[b]{@{}l@{}}
\Axiom$\Gamma_1 \fCenter M : !A \lor !B$
\Axiom$x:!A, \Gamma_2 \fCenter N_1 : !C$
\Axiom$y:!B, \Gamma_3 \fCenter N_2: !C$
\RightLabel{\Elim{\lor}}
\TrinaryInf$\Gamma_1, \Gamma_2, \Gamma_3 \fCenter
\dcase{M}{x}{N_1}{y}{N_2} : !C$
\DisplayProof
\end{tabular}
}
\\[4ex]
\multicolumn{2}{c}{\Axiom$\Gamma \fCenter N: \lfalse$
\RightLabel{\FalseInt}
\UnaryInf$\Gamma \fCenter \abort{!A}{N} : !A$
\DisplayProof}
\end{tabular}
\caption{Règles du système propositionnel intuitionniste
de déduction naturelle étiqueté par des termes \Log{tN2ip}.
Les contextes sont des ensembles finis d'affectations de types
à des variables distinctes. La réunion des contextes dans une conclusion est définie
lorsqu'ils attribuent le même type à chaque variable commune.
Les hypothèses inutilisées peuvent rester dans cette réunion.
Les variables déchargées sont choisies distinctes et absentes
du contexte de conclusion, au besoin par renommage préalable.
L'axiome contextualisé permet les hypothèses inutilisées
et la décharge vacuante.}
\ollabel{tab:tN2ip}
\end{table}

\end{document}
